\documentclass{article}
\usepackage[T1]{fontenc}
\usepackage[utf8]{inputenc}
\usepackage{hyperref}
\usepackage{graphicx}
\hypersetup{pdftitle={Teste 1}}
\title{Teste 1}
\date{}
\begin{document}
\maketitle
\section*{Parte 1}
Primeiro parágrafo desta parte. Este é um parágrafo médio. Lorem ipsum dolor sitamet consectetur adipisicing elit. Est blanditiis incidunt, asperiores errorsuscipit quas odio! Perferendis, commodi. Vel beatae sapiente, sed autemexcepturi temporibus sunt enim ducimus voluptate quaerat eligendi distinctionihil assumenda dolores similique fuga! Possimus, dolorum consectetur?

Outro parágrafo pequeno: Lorem ipsum dolor sit, amet consectetur adipisicingelit. Quibusdam sapiente architecto blanditiis iure ullam hic!

\subsection*{Parte 1.1 <\&>}
Outro parágrafo pequeno: Lorem ipsum dolor sit, amet consectetur adipisicingelit. Quibusdam sapiente architecto blanditiis iure ullam hic!

Aqui fica um link "inline": fica para depois...

Agora vem um link em parágrafo à parte:

Aqui estaria um link. Fica para depois...

A seguir vem a Parte 1.2. Eis como introduzir uma sub-secção em Markdown:

\textbackslash{}\#\# Parte 1.2

(espero que o \textbf{compilador} não tenha \textit{gerado} um h2 )

\subsection*{Parte 1.2}
Aqui começa o 1o parágrafo da Parte 1.2. No código MD, este parágrafo fica bemcoladinho ao título da secção.

\section*{Parte 2: Listas.}
Nesta parte tratamos de listas. Mas isto fica para depois...

\section*{}
O '\#' anterior marca uma secção sem título. No dingus do DaringFireball éapenas um parágrafo. No VS Code e no CommonMark é um h1 vazio. Em geral,nós vamos seguir o Common Mark.\end{document}
